\subsection{Finite-Temperature Phase Behaviour (Task~A)}
\label{subsec:results_phaseA}

Long-time heating and cooling MD cycles were performed on \CsSnI{},
\FASnI{}, and \MASnI{} to characterise finite-temperature structural
evolution.
Three structural observables were tracked as a function of
temperature: the $b/a$ lattice-parameter ratio, the direction-resolved
\ce{Sn-I-Sn} framework linkage distance, and the Glazer-angle
statistics.

Under the common production protocol, \FASnI{} and \MASnI{} remained
dynamically stable through the full heating cycle to \SI{550}{K}
and could be cooled back to their original low-temperature structures
without degradation.
In both hybrid systems, a clear orthorhombic-to-tetragonal transition
was observed at approximately \SI{200}{K} under the present MD
protocol (Table~\ref{tab:phaseA}).
The transition is identified by the simultaneous approach of $b/a$
toward unity and the convergence of the $x$- and $y$-directed
\ce{Sn-I-Sn} linkage distances, which are inequivalent in the
orthorhombic phase and become equivalent in the tetragonal phase.
In \FASnI{}, $b/a$ shifts from 0.985 at \SI{0.1}{K} to 1.000 at
\SI{200}{K}, accompanied by convergence of the $x$- and $y$-directed
linkage distances from
\SI{4.480}{\angstrom}/\SI{4.412}{\angstrom}
to
\SI{4.450}{\angstrom}/\SI{4.449}{\angstrom}
(Table~\ref{tab:fasni_struct}).
A nearly identical pattern is observed in \MASnI{}
(Table~\ref{tab:masni_struct}).
The Glazer-angle standard deviations increase monotonically with
temperature in both systems, reflecting progressive build-up of
dynamic octahedral disorder, while the mean Glazer angles remain
relatively stable below \SI{300}{K}, consistent with the absence of
a clear tetragonal-to-cubic transition under the common protocol.

A higher-temperature tetragonal-to-cubic transition was not clearly
resolved in \FASnI{} or \MASnI{} under the present heating rate
and supercell conditions.
This is consistent with the known sensitivity of simulated
phase-transition temperatures in halide perovskites to heating rate,
supercell size, and force-field details.\cite{fransson2023phase}

\CsSnI{} presented a more demanding case.
The all-inorganic system became dynamically unstable above
approximately \SI{400}{K} under the common protocol, with
progressive loss of the perovskite framework structure.
A separate follow-up using a \ce{Cs}-specific refined model
(without D3) successfully stabilised the simulation to \SI{550}{K},
and the orthorhombic-to-tetragonal transition was recovered in the
\SIrange{200}{300}{K} range (Table~\ref{tab:phaseA};
Figure~S5).
The Glazer-angle standard deviations in the \ce{Cs} follow-up
increase from \ang{0.2} at \SI{0.1}{K} to \ang{10.5} at
\SI{500}{K}, consistent with growing dynamic disorder in the
high-temperature regime.
The greater protocol sensitivity of \CsSnI{} relative to the hybrid
systems likely reflects its higher anharmonicity and more complex
tilt dynamics at elevated temperature, consistent with the anomalously
large thermal expansion and strongly anharmonic lattice behaviour
reported for this compound.\cite{monacelli2023thermodynamics}

\begin{table}[htbp]
  \centering
  \caption{Summary of observed phase-transition behaviour under the
    Task~A MD protocol. Crossover temperatures are workflow-specific
    signatures derived from the present heating/cooling protocol and
    are not calibrated experimental transition temperatures.}
  \label{tab:phaseA}
  \begin{tabular}{llccc}
    \toprule
    System / Run & Model & Stable? &
      Ortho.\ $\to$ Tet.? & Crossover $T$~(K) \\
    \midrule
    \FASnI{} & Common (D3) & Yes & Yes & $\sim$200 \\
    \MASnI{} & Common (D3) & Yes & Yes & $\sim$200 \\
    \CsSnI{} (baseline) & Common (D3) & No & Partial & $\sim$300$^{a}$ \\
    \CsSnI{} (follow-up) & \ce{Cs}-refined (no D3) & Yes & Yes & 200--300 \\
    \bottomrule
  \end{tabular}
  \begin{flushleft}
    \small $^{a}$ Partial transition observed before framework
    degradation at higher temperature.
  \end{flushleft}
\end{table}

\begin{table}[htbp]
  \centering
  \caption{Temperature-dependent structural metrics for \FASnI{}.}
  \label{tab:fasni_struct}
  \begin{tabular}{lcccccc}
    \toprule
    & \multicolumn{6}{c}{Temperature (K)} \\
    \cmidrule(lr){2-7}
    Quantity & 0.1 & 100 & 200 & 300 & 400 & 500 \\
    \midrule
    $\langle b/a \rangle$
      & 0.985 & 1.008 & 1.000 & 1.001 & 1.000 & 0.999 \\
    $\langle d_{\text{Sn-I}}^{x}\rangle$~(\AA)
      & 4.480 & 4.442 & 4.450 & 4.436 & 4.445 & 4.410 \\
    $\langle d_{\text{Sn-I}}^{y}\rangle$~(\AA)
      & 4.412 & 4.449 & 4.449 & 4.439 & 4.445 & 4.417 \\
    $\langle d_{\text{Sn-I}}^{z}\rangle$~(\AA)
      & 6.212 & 6.283 & 6.292 & 6.287 & 6.271 & 6.208 \\
    $\langle \Theta_{\text{Glazer}}\rangle$~($^{\circ}$)$^{a}$
      & 12 & 13 & 13 & 14 & 15 & 16 \\
    $\sigma(\Theta_{\text{Glazer}})$~($^{\circ}$)$^{a}$
      & 6.3 & 6.8 & 6.7 & 7.7 & 8.2 & 9.3 \\
    \bottomrule
  \end{tabular}
  \begin{flushleft}
    \small $^{a}$ Average (and standard deviation) over the three
    Glazer directions $\alpha$, $\beta$, $\gamma$.
  \end{flushleft}
\end{table}

\begin{table}[htbp]
  \centering
  \caption{Temperature-dependent structural metrics for \MASnI{}.}
  \label{tab:masni_struct}
  \begin{tabular}{lcccccc}
    \toprule
    & \multicolumn{6}{c}{Temperature (K)} \\
    \cmidrule(lr){2-7}
    Quantity & 0.1 & 100 & 200 & 300 & 400 & 500 \\
    \midrule
    $\langle b/a \rangle$
      & 1.083 & 0.973 & 1.000 & 1.000 & 1.001 & 1.001 \\
    $\langle d_{\text{Sn-I}}^{x}\rangle$~(\AA)
      & 4.134 & 4.399 & 4.378 & 4.408 & 4.408 & 4.390 \\
    $\langle d_{\text{Sn-I}}^{y}\rangle$~(\AA)
      & 4.476 & 4.278 & 4.375 & 4.406 & 4.406 & 4.390 \\
    $\langle d_{\text{Sn-I}}^{z}\rangle$~(\AA)
      & 6.276 & 6.295 & 6.270 & 6.237 & 6.221 & 6.170 \\
    $\langle \Theta_{\text{Glazer}}\rangle$~($^{\circ}$)$^{a}$
      & 25 & 25 & 21 & 19 & 19 & 20 \\
    $\sigma(\Theta_{\text{Glazer}})$~($^{\circ}$)$^{a}$
      & 4.4 & 6.5 & 8.7 & 9.0 & 9.5 & 10.3 \\
    \bottomrule
  \end{tabular}
  \begin{flushleft}
    \small $^{a}$ Average over $\alpha$, $\beta$, $\gamma$ directions.
  \end{flushleft}
\end{table}

\subsection{Band-Gap Response to Lattice Expansion (Task~B)}
\label{subsec:results_taskB_expansion}

Controlled isotropic volume scaling was applied to all nine
material/phase combinations to isolate the contribution of lattice
size to band gap independent of phase identity or cation orientational
state.
The band gap is plotted as a function of volume-scaling factor for
all materials and phases in Figure~\ref{fig:bg_vs_volume}.

Across all nine cases, lattice expansion consistently opens the band
gap and compression consistently closes it.
This is the clearest and most transferable quantitative relationship
identified in this study.
When evaluated over the \numrange{1.00}{1.10} volume-scaling range,
the band-gap opening rate is broadly similar across most materials
and phases, corresponding to a total increase of approximately
\SIrange{0.07}{0.08}{eV} over a \SI{10}{\percent} volume expansion,
equivalent to a sensitivity of approximately
\SI{0.075}{eV} per \SI{1}{\percent} volume expansion.
A notable deviation is observed for orthorhombic \MASnI{}, which
shows a smaller total response of approximately \SI{0.05}{eV} over
the same \SI{10}{\percent} expansion range.

The positive band-gap response to lattice expansion observed here is
qualitatively consistent with prior literature on \CsSnI{}, which
reports an anomalous increase of band gap with increasing lattice
constant associated with unusually large thermal
expansion.\cite{monacelli2023thermodynamics}
In the framework of crystal-orbital analysis, lattice expansion
reduces \ce{Sn-I} orbital overlap, which raises the antibonding
conduction-band minimum (CBM) and lowers the bonding valence-band
maximum (VBM) in the inverted band ordering characteristic of tin
halides, resulting in gap opening.\cite{huang2013electronic}
Our results extend this qualitative mechanism across the full set of
phases and compositions studied here.

% \begin{figure}[htbp]
%   \centering
%   \includegraphics[width=0.85\linewidth]{figures/bg_vs_volume.pdf}
%   \caption{Band gap as a function of isotropic volume-scaling factor
%     $\alpha$ for all nine material/phase combinations.
%     Filled symbols: orthorhombic; half-filled: tetragonal;
%     open: cubic.
%     Circles: \CsSnI{}; squares: \FASnI{}; triangles: \MASnI{}.}
%   \label{fig:bg_vs_volume}
% \end{figure}

\subsection{Phase Identity Shifts the Band-Gap Baseline (Task~B)}
\label{subsec:results_taskB_phase}

Beyond the common expansion trend, Figure~\ref{fig:bg_vs_volume}
reveals a second important effect: different polymorphs of the same
material do not share the same band-gap baseline at the equilibrium
volume, and the magnitude of the phase-dependent shift is
composition-dependent.

For \CsSnI{}, the phase-to-phase variation in band gap at the
equilibrium volume is comparatively small.
For \FASnI{} and \MASnI{}, the variation is more pronounced.
This difference has a natural structural interpretation: in hybrid
tin perovskites, phase transitions involve not only changes in
octahedral tilt pattern but also reorganisation of the A-site
cation environment, producing larger changes in local \ce{Sn-I}
geometry --- and hence larger electronic-structure shifts --- than
in the all-inorganic case.

Critically, this finding means that the temperature-dependent
band-gap evolution cannot be accounted for by lattice expansion
alone, even if the thermal expansion coefficient were precisely
known.
As the material heats through a phase transition, the band-gap
baseline shifts discontinuously to that of the new polymorph, and
this shift can either reinforce or partially offset the gradual
expansion-driven opening, depending on whether the new-phase
baseline is higher or lower than that of the preceding phase.
A first-order understanding of temperature-dependent band-gap
behaviour in these materials therefore requires knowledge of both
the expansion sensitivity and the phase-dependent baseline offsets.

\subsection{Organic-Cation Orientation and Band Gap (Task~C)}
\label{subsec:results_taskC}

\subsubsection{Tetragonal-Phase Study}

The sensitivity of the band gap to organic-cation orientation was
investigated in the tetragonal phase of \FASnI{} and \MASnI{}
using the synchronised rotation protocol described in
Section~\ref{subsec:taskC}.
Results are shown in Figure~\ref{fig:taskC_tetra}.

\FASnI{} shows relatively weak orientational sensitivity in the
tetragonal phase.
The band gap remains near \SI{0.5}{eV} across the full ordered
rotation range, indicating that the FA cation does not strongly
perturb the \ce{Sn-I} cage through ordered in-plane rotation in
this phase.
The randomly rotated FA reference gives a band gap of approximately
\SI{0.58}{eV}, slightly above the ordered-rotation trend, suggesting
that disordered FA orientations produce a somewhat larger electronic
perturbation than ordered ones under the tetragonal symmetry
constraint.

\MASnI{} behaves very differently.
In the ordered-rotation tetragonal study, the band gap spans
approximately \SIrange{0.18}{0.75}{eV} as a function of rotation
angle, demonstrating that the MA cation couples far more strongly
to the inorganic cage under the synchronised-rotation assumption.
The randomly rotated MA reference gives a band gap of approximately
\SI{0.62}{eV}, which falls within the range of the ordered-rotation
scan.
The large span of the MA ordered-rotation curve is consistent with
the smaller and more aspherical geometry of the MA cation relative
to FA, which may allow stronger orientation-specific hydrogen
bonding and cage distortion at particular rotation angles.
\cite{maheshwari2019relation}

It is important to note that the synchronised-rotation protocol
is more directly motivated for the FA case, where ordered
low-temperature arrangements are experimentally documented, than for
MA, where the ordered pathway is less definitively established.
\cite{schueller2018crystal}
The tetragonal MA result should therefore be interpreted as a
controlled hypothesis test of orientational sensitivity rather than
a confirmed physical rotation pathway.

% \begin{figure}[htbp]
%   \centering
%   \includegraphics[width=0.85\linewidth]{figures/taskC_tetragonal.pdf}
%   \caption{Band gap as a function of synchronised cation rotation
%     angle for \FASnI{} (blue) and \MASnI{} (red) in the tetragonal
%     phase.
%     Filled circles: ordered rotation scan.
%     Stars: randomly rotated reference.}
%   \label{fig:taskC_tetra}
% \end{figure}

\subsubsection{High-Temperature Cubic-Phase Study}

In the high-temperature cubic phase, a mean-angle sweep over
20 orientation configurations was performed for both materials.
Strikingly, the relative orientational sensitivities of MA and FA
reverse compared to the tetragonal case
(Figure~\ref{fig:taskC_cubic}).

In the cubic phase, \MASnI{} shows only weak sensitivity: the band
gap spans approximately \SIrange{0.68}{0.76}{eV}, a total variation
of approximately \SI{0.076}{eV}.
By contrast, \FASnI{} shows a significantly larger response, with
band gaps spanning approximately \SIrange{0.50}{0.67}{eV}, a total
variation of approximately \SI{0.17}{eV} --- more than double the
MA range under comparable orientation conditions.

This reversal of the MA/FA sensitivity ordering between the
tetragonal and cubic regimes is a central result of this study.
It implies that the coupling between cation orientation and band gap
is not a fixed property of the cation species but depends critically
on the structural context provided by the surrounding phase.
A plausible interpretation is that the larger and more deformable
FA cation couples more effectively to the cage in the high-symmetry
cubic environment, where the tilt constraints of the lower-symmetry
phases are relieved and the cage cavity geometry allows stronger
orientation-dependent interactions.
Conversely, in the tetragonal phase, the constrained tilt pattern
may preferentially amplify the coupling of the smaller MA cation to
the cage through specific hydrogen-bond-driven distortions.
While this interpretation is consistent with the known differences
in cation geometry and rotational dynamics between MA and
FA,\cite{laven2023rotational, maheshwari2019relation} the precise microscopic
origin of the sensitivity reversal remains to be established and
is the subject of future work.

% \begin{figure}[htbp]
%   \centering
%   \includegraphics[width=0.85\linewidth]{figures/taskC_cubic.pdf}
%   \caption{Band gap as a function of mean cation orientation angle
%     for \FASnI{} (blue circles) and \MASnI{} (red squares) in the
%     high-temperature cubic phase.
%     The $x$-axis parameterises the 20 mean-angle configurations
%     used in the sweep (see Section~\ref{subsec:taskC}).}
%   \label{fig:taskC_cubic}
% \end{figure}

\subsection{Integrated Temperature-and-Strain Response (Task~D)}
\label{subsec:results_taskD}

Task~D combined finite-temperature structural sampling with
externally imposed isotropic strain to provide a first view of
band-gap behaviour under conditions that simultaneously include
thermal disorder, phase identity, and applied strain.
Table~\ref{tab:taskD} summarises the band-gap values obtained for
all materials across three phase-representative temperatures and
three strain states.

\begin{table}[htbp]
  \centering
  \caption{Band gap (eV) as a function of material, phase,
    temperature, and imposed isotropic strain ($\alpha$).
    Values represent the mean of two sampled SS-DFT structures
    per condition.}
  \label{tab:taskD}
  \begin{tabular}{llcSSS}
    \toprule
    Material & Phase & $T$~(K) &
      {$\alpha=0.95$} & {$\alpha=1.00$} & {$\alpha=1.05$} \\
    \midrule
    \CsSnI{} & Cubic        & 400 & 1.684 & 1.722 & 1.772 \\
    \CsSnI{} & Tetragonal   & 325 & 0.906 & 0.937 & 0.951 \\
    \CsSnI{} & Orthorhombic & 250 & 0.268 & 0.397 & 0.523 \\[4pt]
    \FASnI{} & Cubic        & 320 & 0.755 & 0.825 & 0.800 \\
    \FASnI{} & Tetragonal   & 250 & 0.415 & 0.819 & 0.441 \\
    \FASnI{} & Orthorhombic & 150 & 0.577 & 0.704 & 0.837 \\[4pt]
    \MASnI{} & Cubic        & 320 & 0.701 & 0.853 & 0.755 \\
    \MASnI{} & Tetragonal   & 250 & 0.785 & 0.691 & 0.733 \\
    \MASnI{} & Orthorhombic & 150 & 0.718 & 0.608 & 0.761 \\
    \bottomrule
  \end{tabular}
\end{table}

\subsubsection{\CsSnI{}}

\CsSnI{} exhibits the most regular behaviour in Task~D.
The unstrained band gap increases monotonically across the three
phase-representative conditions: \SI{0.397}{eV} (orthorhombic,
\SI{250}{K}), \SI{0.937}{eV} (tetragonal, \SI{325}{K}), and
\SI{1.722}{eV} (cubic, \SI{400}{K}).
At each temperature, the strain response is monotonic and consistent
with the Task~B result: tensile strain gives the highest band gap
and compressive strain the lowest ($\alpha=0.95 < \alpha=1.00 <
\alpha=1.05$ at all three temperatures).
The all-inorganic system is thus well described by the simple picture
in which lattice expansion controls the first-order band-gap trend
and phase identity determines the baseline from which expansion acts.

\subsubsection{\FASnI{}}

\FASnI{} shows a more complex response than \CsSnI{}.
While the unstrained band gap increases overall with temperature
(\SI{0.704}{eV} at \SI{150}{K}, \SI{0.819}{eV} at \SI{250}{K},
\SI{0.825}{eV} at \SI{320}{K}), the strain dependence is not
uniform across phase conditions.
In the orthorhombic case at \SI{150}{K}, the ordering follows the
expected trend (compressive < unstrained < tensile).
In the tetragonal case at \SI{250}{K}, the response is strongly
anomalous: the unstrained structure gives the \emph{highest} band
gap (\SI{0.819}{eV}), while both compressive (\SI{0.415}{eV}) and
tensile (\SI{0.441}{eV}) strain reduce it substantially.
In the cubic case at \SI{320}{K}, the ordering returns closer to
the expected form, although the tensile value remains slightly below
the unstrained reference.

\subsubsection{\MASnI{}}

\MASnI{} shows the largest qualitative deviation from a simple strain
picture.
Not only does the sign of the anomaly change across phase conditions,
but the direction differs from that in \FASnI{}.
At \SI{250}{K} (tetragonal), compressive strain ($\alpha=0.95$)
gives the \emph{highest} band gap (\SI{0.785}{eV}), while the
unstrained structure gives the lowest (\SI{0.691}{eV}).
At \SI{150}{K} (orthorhombic), the unstrained structure gives the
lowest value and tensile strain the highest, consistent with
Task~B, but the differences are anomalously small.
At \SI{320}{K} (cubic), the unstrained structure gives the highest
band gap, while compressive strain gives the lowest, consistent with
the Task~B expectation but with an inverted tetragonal result.

\subsubsection{Mechanistic Interpretation of Task~D}

The divergent behaviour of the hybrid systems relative to \CsSnI{}
is naturally interpreted in light of the Task~C results.
Unlike the idealised lattice-scaling study in Task~B, the Task~D
protocol samples finite-temperature configurations that include not
only the average lattice dimensions set by the imposed strain, but
also additional structural variance from thermal disorder ---
including local octahedral distortions and, in the hybrid systems,
partially disordered cation orientations.
The Task~C results showed that cation orientational effects can
shift the band gap by \SIrange{0.076}{0.17}{eV}, which is
comparable in magnitude to the strain effects observed here.
When thermal cation disorder is present, the band-gap response to
imposed strain therefore reflects the combined action of lattice
deformation and orientational coupling, and these contributions can
either reinforce or partially cancel depending on the specific
structural configuration sampled.

These results carry an important practical implication.
Manufacturing-induced stresses --- arising from thermal cycling
during module lamination, lattice mismatch at interfaces, or
differential thermal expansion with adjacent layers --- may perturb
the band gap of tin-halide perovskite absorbers by amounts that are
not predictable from the zero-temperature strain response alone.
In the hybrid systems, the magnitude and direction of the
perturbation will depend on the thermal history through its coupling
to cation disorder, suggesting that controlling cation dynamics ---
through A-site engineering or processing protocols that influence
cation ordering --- may offer an additional lever for mitigating
undesirable strain-induced band-gap shifts in device fabrication.
\cite{yang2022strain}

We note that the Task~D results are based on a limited number of
sampled structures (two per condition) and should be treated as
preliminary finite-temperature observations rather than fully
converged ensemble averages.
Statistical convergence of the hybrid-system behaviour in particular
requires larger ensembles and is the subject of ongoing work.

\subsection{Composite Mechanism of Temperature-Dependent Band-Gap Evolution}
\label{subsec:composite}

Taken together, the four tasks converge on a single overarching
conclusion: the net temperature dependence of band gap in tin-halide
perovskites is a composite effect that cannot be reduced to any
single structural mechanism.
Figure~\ref{fig:composite_schematic} illustrates this picture
schematically.

Starting from the low-temperature equilibrium structure, heating
simultaneously drives:
(i) an increase in unit cell volume through thermal expansion,
which opens the band gap at a rate of approximately
\SI{0.075}{eV} per \SI{1}{\percent} volume expansion;
(ii) a phase transition, which shifts the band-gap baseline by an
amount that is composition-dependent and can either reinforce or
partially offset the expansion-driven opening; and
(iii) in hybrid systems, an increase in cation orientational
disorder that can shift the band gap by up to \SI{0.17}{eV}
depending on cation species and phase regime (Section~\ref{subsec:results_taskC}).

For \CsSnI{}, the absence of organic-cation effects reduces the
composite behaviour to the simpler combination of expansion and
phase identity, producing the comparatively regular temperature
evolution observed in Tasks~B and D.
For \FASnI{} and \MASnI{}, the cation disorder contribution adds
an independent term that becomes particularly important in the
high-temperature cubic phase, where orientational sensitivity is
largest (Task~C).
The MA/FA contrast --- with MA showing stronger orientational
coupling in the tetragonal phase and FA showing it in the cubic
phase --- means that the two hybrid systems are not interchangeable
from the perspective of thermal band-gap management, even though
they share the same basic perovskite framework.

This composite picture has a direct implication for experimental
measurements of temperature-dependent band gap: observed trends that
appear to follow a single smooth function of temperature may in fact
reflect the accidental cancellation or reinforcement of multiple
contributions with very different functional forms.
Disentangling these contributions experimentally requires
measurements that separately resolve structural evolution (phase
identity, lattice parameters) and electronic structure --- for
example, variable-temperature diffraction combined with optical or
photoemission spectroscopy.

% \begin{figure}[htbp]
%   \centering
%   \includegraphics[width=0.90\linewidth]{figures/composite_schematic.pdf}
%   \caption{Schematic summary of the three structural contributions
%     to temperature-dependent band-gap evolution in tin-halide
%     perovskites studied in this work.
%     The net band gap at any temperature reflects the superposition
%     of (i) the expansion-driven opening trend, (ii) the
%     phase-dependent baseline shift at each transition, and
%     (iii) in hybrid systems, the cation-disorder perturbation.
%     The relative magnitudes and signs of contributions (ii) and (iii)
%     are material-specific, which explains why different perovskite
%     compositions can exhibit qualitatively different
%     temperature-dependent band-gap behaviour.}
%   \label{fig:composite_schematic}
% \end{figure}